\section{部署与信任}

\subsection{部署复盘与回滚}

每次 release 后都需要做 pre-mortem、post-mortem 与 rollback 验证：
\begin{itemize}
    \item \textbf{Pre-mortem}：在 rollout 之前模拟 worst-case prompt drift，确认 guardrail 覆盖。
    \item \textbf{Post-mortem}：收集 rejection log、reward score drift，在 incident log 中备注。
    \item \textbf{Rollback 验证}：设计 rollback playbook，先在 staging 里演练 restore checkpoint、recover guardrail。
\end{itemize}

\begin{importantbox}{预拟事故的价值}
在产品经理批准 rollout 之前先做一次预估风险演练，可以把 high-risk prompt 列入 watchlist，避免事故发生时手忙脚乱。
\end{importantbox}

\begin{table}[H]
\centering
\begin{tabular}{p{0.25\textwidth} p{0.65\textwidth}}
\toprule
复盘阶段 & 内容 \\
\midrule
Pre-mortem & drift risk list、evidence matrix 检查、guardrail traceability \\
Post-mortem & reward score drift graph、human override log、root cause analysis \\
Rollback & reference checkpoint + runbook、alert escalation path \\
\bottomrule
\end{tabular}
\caption{部署复盘与回滚交付物}
\end{table}

\subsection{多角色治理与信任}

Alignment pipeline 不只是 research，治理团队、法律、SRE 都需要参与：
\begin{itemize}
    \item 研发团队负责 reward model + policy update。
    \item 运营/监控团队负责 drift alert、dashboard。
    \item 法务/合规负责 guardrail audit、privacy review。
    \item PM + decision board 根据 evidence ranking 最终决定 rollout 或 rollback。
\end{itemize}

\begin{knowledgebox}{多角色治理的设想}
把 evidence matrix 声明为 multi-stakeholder checkpoint，例如：每次 policy update 都要让 compliance team 复核 long-form prompts。
\end{knowledgebox}

\begin{warningbox}{风险放大链条}
如果缺乏 governance checkpoint，就会出现 reward model fine-tune 后拉高某类 prompt，但 compliance 没有 catch，导致 high-risk reply 直接进入 production。
\end{warningbox}

\subsection{本章小结}

部署与信任需要 pre-mortem/post-mortem+rollback 以及多角色的 evidence governance，只有这样才能在不断迭代的 RLHF pipeline 中维持 alignment 信任。

